\cotefile{39}
\scanpage{3}
% 1
$k$ un \deleted{corps \inserted{\deleted{anneau}} anneau avec} corps

% 2–3
$S_n(k)=k[[t_1,\ldots,t_n]]$\qquad $(n\geq0\ ;\ S_0(k)=k)$\qquad pour $n$ variable, on a anneau cosimplicial

% 4
$A_n\subset S_n(k)$\qquad sous-anneau cosimplicial i.e.

% 5
\margin{$\Big\downarrow$}% curved relocation arrow from (b) to the space after the three clauses
(b) $\forall n\geq0$, $A_n$ est un sous anneau unifère de $S_n$

% 6–9
(a) On a\quad $a_1)$ $A_n$ est stable par $\unsure{\mathfrak S|\mathscr S}_n$ opérant sur $\unsure{S_n|\text{les}}$

\hspace*{2em}$a_2)$ Si $f(t_1,\ldots,t_n)\in A_n$, alors $f(\underbrace{t_1,\ldots,t_n}_{\text{\inserted{$t_{n+1}$ variable qui ne figure pas}}})\in A_{n+1}$

\hspace*{2em}$a_3)$ Si $f(t_1,\ldots,t_n)\in A_n$, alors $f(t_1,t_2,\ldots,\underbrace{t_{n-1},t_{n-1}}_{\text{répété}})\in A_{n-1}$

% 10–11
Le plus petit $A_*=(A_n)_{n\in\mathbb N}$ \deleted{\unsure{pareil|pareille|famille}} \inserted{satisfaisant a) b)} est donné par \deleted{$A_n=\unsure{\mathbb P|\mathbb F}$}
\[
 A_n=\mathrm{Im}(\mathbb Z\to S_n)
\]

% 12–13
donc l'\deleted{objet}\inserted{anneau} cosimplicial constant de valeur $\mathrm{Im}(\mathbb Z\to k)$ (qui est $\mathbb Z$ ou un $\mathbb Z/p\mathbb Z$, \deleted{\illegible{2}}\inserted{$p$ nb premier} $=\mathrm{car}.\,k$)

% 14–17
On supposera de plus que

c)\qquad Si $f\in A_n$, \inserted{et} $f\in\mathrm{rad}\, S_n$, i.e. \deleted{\unsure{est}} $f$ inv. dans $S_n$, alors $\frac1f\in S_n$

$\Updownarrow$

c')\qquad $A_n$ est un anneau local, et l'inclusion $A_n\to S_n$ est un hom local

% 18–21
Le plus petit $A_*$ satisfaisant a) b) c) est formé des
\[
 A_n=\mathrm{Im}(\mathbb P\to S_n)
\]
où $P$ est le sous-corps premier de $k$ ($\mathbb Q$ en car $0$, $\mathbb F_p$ en car $p>0$).

% 22
d)\qquad $t_1\in A_1\subset S[[t_1]]$

% 23–24
Le plus petit $A_*$ satisfaisant à a) b) d) \deleted{\unsure{est|c)}} est formé des
\[
 A_n=\mathbb P[t_1,\ldots,t_n]_{\deleted{\unsure{\mathfrak m_n}}}\qquad
 \deleted{\unsure{\mathfrak m_n=(}}
\]

% 25–26
Le plus petit $A$ satisfaisant a) b) c) d) est formé des
\[
 A_n=\mathbb P[t_1,\ldots,t_n]_{\mathfrak m_n}\qquad
 \mathfrak m_n=\sum t_i\,\mathbb P[t_1,\ldots,t_n].
\]

% 27–30
\deleted{\unsure{a)|e)}} Un hom $S_n\longrightarrow S_m$ \inserted{ou $\underset{\unsure{\text{``}D_m\text{''}|D'_m|X_m}}{\mathrm{Specf}(S_m)}-\underset{\unsure{\text{``}D_n\text{''}|D''_n|X_n}}{\mathrm{Specf}(S_n)}$} donné par $f_1,\ldots,f_n\in S_m$ est dit un $A$-hom si $f_1,\ldots,f_n\in A_m$. Il résulte de b) et d) que l'identité $D_n\to D_n$, \deleted{est un \unsure{et les}}\inserted{\deleted{\unsure{tout applicati.|toute application}}plus gén. les applications} $D_m\to D_n$ sont \deleted{\unsure{des|donc}}\inserted{des} $A$-hom ($\varepsilon_i\in A_n$ pour $1\leq i\leq n$). L'analogue de b)
% archive sheet number: 2
\scanpage{4}
% 1–3
est la condition

e) Le composé de deux \inserted{$A_*$-}morphismes $X_m-X_n-X_p$ est un $A_*$-morphisme.

% 3, left
% slanted
\margin{Moyennant a) à e) $\mathrm{Hom}_A(X_m,X_n)\ \unsure{\simeq|\approx}\ \mathrm{Hom}_{A_0\text{-alg}}(A_n,A_m)$ \unsure{?}}

% 4–8
\uline{NB} Moyennant \inserted{b) et} e), la condition a) équivaut au fait que
\[
\left\{\begin{aligned}
 a_1)&\quad 1\in A_0\\
 a_2)&\quad \text{le morphisme ``différence}\quad X_1\times X_1-X_1\text{ est } A_*\text{-morph.}\\
 a_3)&\quad \text{--------------- ``produit''}\quad X_1\times X_1-X_1\text{ est } A_*\text{-morph.}
\end{aligned}\right.
\]

% 9–10
\uline{NB} Les $A_n=k[t_1,\ldots,t_n]$ ou $A_n=k[t_1,\ldots,t_n]_{\mathfrak m_n}$ satisfont e)

% 11–15
L'analogue de c) est la condition

f) Si $u:X_n\longrightarrow X_n$ est un $A$-morphisme et un isomorphisme, \deleted{alors} i.e. le déterminant \deleted{\unsure{fct|fon}} de la matrice jacobienne de $u$ est inversible, alors $u^{-1}$ est un $A$-morphisme

% 16–17
Cette condition implique \deleted{(\unsure{et est|et est ce} équivalente à ?)}

% 16, left
% slanted
\margin{f') implique-t-il f (moyennant a) à e)) ? il ne semble pas}

% 18–20
f') Les $A_n$ sont des anneaux henséliens

[En fait, si f) est satisfait pour un $n$, (f') est satisfait pour cet $n$, i.e. $A_n$, est hensélien.]

% 20–24
\deleted{D'ailleurs si $u:\inserted{X_n\xrightarrow{\sim}X_n}$ est donné par $f_1,\ldots,f_n\in B_n$, alors $u^{-1}$ est donné par $g_1,\ldots,g_n$ appartenant à la clôture hensélienne de $(k[f_1,\ldots,f_n]_{\mathfrak n}$, où $\mathfrak n=k[f_1,\ldots,f_n]$.}

% 22, margin
% slanted
\margin{\deleted{\unsure{vérifier}}}

% 25–28
\uline{NB} Le plus petit $A_*$ satisfaisant a) à d) et f') est formé des
\[
 B_n=\mathbb P\{t_1,\ldots,t_n\}\qquad
 \left(\begin{aligned}&\text{clôture hensélienne de}\\[-2pt]&\mathbb P[t_1,\ldots,t_n]_{\mathfrak m_n}\end{aligned}\right)
\]
et il satisfait $=$ e) et f).

% 29
\deleted{g) Soit $\xi$ un champ de vecteurs \unsure{for|fon}}

% 30–33
Sous a) b) c) d) e) on définit la catégorie des $A_*$-schémas formels\inserted{et des Modules dessus, fibrés de t.f. !}, et le calcul différentiel dessus (opérateurs différentiels, parties principales, \inserted{fibré tangent} etc ....)

% 31, left
% slanted
\margin{\unsure{(des graphismes)|des graphismes)|6 des propriétés} \unsure{Th.|Obs.} des produits fibrés transversaux etc ....}
% archive sheet number: 3
\scanpage{5}
% 1–6
Nous supposons maintenant $k$ de car. $0$, de sorte qu'on a la notion d'exponentielle d'un champ de vecteurs formel $\xi$ sur \deleted{\unsure{X_n|X}} un schéma formel lisse $X$,
\[
 (t,x)\mapsto(\exp t\xi).x\ :\ \underbrace{X_1}_{\text{``axe temps}}\times X\longrightarrow X
\]

% 7, left
% slanted
\margin{Intégration d'équations différentielles}

% 7–10
g) Si $X$ est un $A$-schéma formel (lisse) et $X$ un \inserted{$A$-}champ de vecteurs, alors la résolvante $X_1\times X\longrightarrow X$ est un $A$-morphisme.

% 11–12
(Bien entendu, il suffirait de le \deleted{faire}\inserted{demander} pour $X$ de la forme $X_n$\ldots).

% 13–16
Ceci implique \inserted{\unsure{et|est}} (par l'intro$\unsure{^{\mathrm n}}$ de la variable temps, devenant variable spatiale), \unsure{donc|son} équivaut, à

g') \unsure{Idem que|Comme|Analogue à} g), mais avec \deleted{``}un ``champ de vecteurs variable avec le temps''

% 17–20
$\xi(t)$\quad (
\diagram{\begin{tikzcd}[ampersand replacement=\&,column sep=small,row sep=small]
X_1\times X \arrow[rr] \arrow[dr,dashed,"\mathrm{pr}_2"'] \& \& \underbrace{\unsure{\mathbb T|T|\mathbf T}(X)}_{\substack{\text{variété }\deleted{\text{\unsure{actuelle|vectorielle|formelle}}}\text{\inserted{formelle}}\\\text{tangente de }X}} \arrow[dl,dashed] \\
\& X \&
\end{tikzcd}}
), en résolvant l'équation différentielle formelle
\[
 \frac{d}{dt}x(t)=\xi(t,x)
\]

% 21–23
La résolvante $x=\varphi(t,x_0)$ correspond à un $\varphi:X_1\times X\longrightarrow X$ qui est un $A$-morphisme.
% archive sheet number: 4
\scanpage{6}
% 1–3
\uline{Conséquence}. Soit $k_0=A_0$, c'est un sous-corps de $k$. Soit $L(k_0)$ la catégorie des Alg. de Lie de dim finie sur $k_0$. Par Hausdorff on a un foncteur

% 4
\diagram{\begin{tikzcd}[ampersand replacement=\&,column sep=large]
L(k_0)\arrow[r,"\text{Hausdorff}"'] \& \mathrm{Grf}(k_0)\arrow[r,"\substack{\text{ext. de la}\\\text{base}}"'] \& \mathrm{Grf}(k)
\end{tikzcd}}

% 5–7
où ``Grf'' désigne la catégorie des groupes formels. Ceci dit, ce foncteur se factorise canoniquement par les $A$-groupes (formels) sur $k$.

% 8–10
De façon plus précise, soit $\overline W(\mathfrak g_0)$ le complété à l'origine \deleted{d'un}\inserted{de $\mathfrak g_0\otimes_{k_0}k$} pour \inserted{un} $k_0$-espace vectoriel \inserted{de rg fini} \deleted{$\unsure{k_0|\mathfrak g_0}$}. C'est un $A$-schéma formel.

% 11–14
Les formules de Hausdorff, si $\mathfrak g_0$ est une Alg de Lie, définissent une loi de groupe
\[
 \overline W(\mathfrak g_0)\times\overline W(\mathfrak g_0)\longrightarrow\overline W(\mathfrak g_0)
\]
\uline{Cette dernière est une $A$-loi} ; \deleted{Pour ceci, il suffit de vérifier \deleted{deux choses}\inserted{que le système de Pfaff}
1°) associé est donné par une $A$-loi \inserted{\unsure{un $A$-fibré}}} et si

% 15–20
$\mathfrak g_0\longrightarrow\mathfrak h_0$ est un hom d'algèbres de Lie sur $k_0$, l'hom de groupes formels $\overline W(\mathfrak g_0)\longrightarrow\overline W(\mathfrak h_0)$ est un $A$-hom. Utilisant le Th. d'Ado, on est ramené à prouver que si $\mathfrak g_0\subset\mathfrak{gl}(n,k_0)$, alors le Noyau de \inserted{ss/}gpe formel de $Gl(n,k_0)$ qu'il définit est un sous-$A$-schéma formel ....
% archive sheet number: 5
\scanpage{7}
% 1–3
\uline{Cor} $A_1\subset S_1(k)$ stable par primitivation canonique (i.e. celle qui s'annule à l'origine), contient $\log(1+t_1)$, $\exp t_1$, $\cos t_1$, $\sin t_1$, $\mathrm{tg}\,t_1$ ....

% 4–8
Soit $X$ schéma formel lisse de dim $n=p+q$. On sait \deleted{que} ce que c'est qu'un \uline{feuilletage} \inserted{de dim $p$} de $X$ : variété formelle quotient $B$ de $X$, avec $X-B$ submersif $\}$, $\dim B=q=n-p$. A un tel feuilletage est associé un sous-module $D\subset\Theta(X)$

% 9–12
[champs de vecteurs le long de la fibre] qui est un ``système de Pfaff de dim $p$ complètement intégrable'', et vice-versa. Si $X$ a une $A$-structure, on dit qu'on a un $A$-feuilletage si $\exists$ \inserted{sur} $B$ une $A$-structure

% 13–17
(d'ailleurs unique \unsure{$\}$|)|,} grâce à f)) telle que $X\to B$ soit un $A$-morphisme. Cela \uline{implique} que le sous-module $D\subset\Theta(X)$ (ou sous-fibré de $T(X)$) soit un $A$-sous-fibré -- i.e. qu'on ait un $A$-syst. de Pfaff. compl. intégrable. On veut qu'on ait une réciproque, savoir

% 18–20
h) Si $D_A$ est un $A$-syst. de Pfaff compl. intégrable sur le $A$-schéma formel $X$, alors il donne lieu à un $A$-feuilletage \deleted{\unsure{complèt$^{\mathrm{t}}$ intégr.|complét. intégr.|réalisé}}. [En d'autres termes,

% 18, left
% slanted
\margin{\uline{NB} h) implique g), qui équivaut à h) pour $p=1$}

% 21–27
si $i:Y-X$ est un $A$-morphisme donné tel que le composé
\[
 T(Y)_0\mathrel{\overset{T(i)_0}{\text{---}}}T(X)_0-T(X)_0/D(0)
\]
soit un isom, alors l'unique morphisme
\[
 \pi:X-Y
\]
tel que $\pi\circ i=\mathrm{id}_Y$ et que $\varphi$ soit compatible avec $D$, est un $A$-morphisme ....

% 28, foot
\margin{$\Longleftarrow$} h') la fibre de $0$ dans un tel feuilletage est une sous-$A$-variété formelle ....]
% archive sheet number: 6
\scanpage{8}
% 1–4
\deleted{\uline{Cor} $A_1$ est stable par primitives canoniques (i.d. prises d'augmentation nulle) contient $\log(1+t_1)$\unsure{,|;} (primitive de $\frac{1}{1+t_1}$) donc $\exp.t_1$ (série inverse), les fonctions $\cos t_1$, $\sin t_1$, $\mathrm{tg}\,t_1$ sont $\in A_1$ .....}

% 5–6
\deleted{\uline{Lemme} \deleted{\unsure{h)|g)|1)} Soit} Soit $X$ un \inserted{$A$-}schéma formel (lisse) \inserted{de dim $n=p+q$}, $D$ un ``système de Pfaff de degré $p$ complètement intégrable sur $X$'', \deleted{\unsure{D}}}

% 7–10
\deleted{\deleted{deux \unsure{soient|lisses} Soient $Y,Z$ des \inserted{$A$-}schémas formels \inserted{(lisses)} de dims $p$ resp $q$,} $p:X\to Y$ et $i:Z\to X$ \inserted{des $A$-morphismes}, tels que $T(p)(0)$ soit un \deleted{\unsure{appl.}}\inserted{isom} $D(0)-T(Y)_0$ et que \deleted{$T(i)(0)$} le composé \deleted{$T(i)(0)$}}
\[
 \deleted{T(Z)_{0}\deleted{(0)}\xrightarrow{T(i)_0}T(X)_{0}\deleted{(0)}\xrightarrow{\mathrm{can}}T(X)_0/D(0)}
\]

% 11–14
\deleted{soit un isom. On sait qu'il existe alors un isomorphisme \inserted{formel (pas néc. $A$-isom)} unique}
\[
 \deleted{X\overset{\varphi}{\simeq}Y\times Z}
\]
\deleted{\deleted{\unsure{transformant|transportant} et \unsure{qui|puis} tel que $\mathrm{pr}_1\circ\varphi=p$, $\varphi\circ i=(z\mapsto(0,z))$} et}

% 15–19
\deleted{Soit $Z$ \inserted{$A$-}schéma formel \inserted{(lisse)} de dim $q$, $i:Z-X$ une $A$-\unsure{immersion|réimmersion} telle que le composé}
\[
 \deleted{T(Z)_0\xrightarrow{T(i)_0}T(X)_0\longrightarrow T(X)_0/D(0)}
\]
\deleted{soit un isom. On sait qu'il existe alors un unique morphisme formel}
\[
 \deleted{X\xrightarrow{\pi}Z}
\]

% 20–25
\deleted{tel que $\pi\circ i=\mathrm{id}$, et que ``$\pi$ soit \deleted{\unsure{constante sur|constant sur|constant sur les}}\inserted{compatible avec} le feuilletage défini par $D$''. Conditions équivalentes}

\deleted{(i) $\pi$ est un $A$-morphisme}

\deleted{(ii) $\exists$ $A$-isom $X\simeq B\times Y$, tel que le feuilletage formel de $X$ défini par $D$ soit le feuilletage défini par la projection $\mathrm{pr}_1:X\to B$. On dit}
% archive sheet number: 7
\scanpage{10}
% 1–5
$k$ anneau commutatif
\[
\begin{aligned}
 S_n&=k[[t_1,\ldots,t_n]]\qquad(n\geq0)\\
 X_n&=\mathrm{Specf}(S_n)\\
 X_*&=(X_n)_{n\geq0}\quad\text{schéma formel simplicial augmenté}\\
 S_*&=(S_n)_{n\geq0}\quad\text{anneau top. cosimplicial coaugmenté}
\end{aligned}
\]

% 6–7
On suppose donné, pour tout $n\inserted{\geq0}$, un sous-ensemble
\[
 A_n\subset S_n
\]

% 8–11
On pose $A_*=(A_n)_{n\geq0}$. Un morphisme $X_m\xrightarrow{p}X_n$, associé à un hom de $k$-algèbres topologiques $S_n\xrightarrow{p^*}S_m$, est dit un $A_*$-\uline{morphisme} si $p^*(t_i)=f_i\in S_m$ est $\in A_m$ pour tout $1\leq i\leq n$.

% 11, left
% slanted
\margin{\uline{NB} les $f_i$ sont tels que $\varepsilon(f_i)$ \uline{nilpotents}}

% 12–13
Donc cela signifie que les $\mathrm{pr}_i\circ p$ $(1\leq i\leq n)$ sont des $A_*$-morph.

On suppose :

% 14–17
% slanted
\margin{Identités}
(a) Les opérations simpliciales entre les $X_n$ sont des $A_*$-morphismes ($\Longleftrightarrow\ \forall n\geq0$ et $1\leq i\leq n$ on a $t_i\in A_n$ ; $\Longleftrightarrow$ les morphismes identiques des $X_n$ sont $\in A_n$)

% 18–22
% slanted
\margin{Composition}
(b) Le composé de deux $A_*$-morphismes est \unsure{itou} [$\Longleftrightarrow$ si $p:X_m\to X_n$ est un $A_*$-morphisme, alors $p^*(A_n)\subset A_m\ \Longleftrightarrow$ Si $f_1,\ldots,f_n\in A_m$, et $\varepsilon(f_1),\ldots,\varepsilon(f_n)$ nilpotents, alors pour tout $\varphi\in A_n$ \inserted{avec $\varepsilon(\varphi)$ nilpotent}, on a $\varphi(f_1,\ldots,f_n)\in A_m$]

% 22, left
% slanted
\margin{NB \deleted{\unsure{d'ailleurs}}\inserted{Ce} \deleted{\unsure{que ce soit|qui est}}\inserted{sera} vrai \inserted{grâce à b)}, (sans supposer $\varepsilon(\varphi)$ nilpotent \deleted{dès que}) dès que $\mathrm{Im}(A_0\to S_n)\subset A_n$ \deleted{si $a\in A_0$, on}}

% 23–29
% slanted
\margin{\unsure{Condition annelée|Condition annelé}}
(c) \inserted{$1\in\unsure{A_0|A_n}$, et} \deleted{Le morphisme ``différence'' $X_2=X_1\mathbin{\widehat{\times}}_k X_1\to X_1$ $((x,y)\mapsto x-y)$ et ``produit'' [$(x,y)\mapsto xy$] sont des $A_*$ morphismes [\,$\Leftrightarrow t_1-t_2$ et $t_1t_2$ sont $\in A_2$, \inserted{et $1\in A_0$} $\Longleftrightarrow$} \deleted{les}\inserted{les} $A_n$ sont \deleted{stables par \deleted{\unsure{sommes et}} \inserted{différences et} produits, \deleted{et par différences} \inserted{et}}\inserted{des} \deleted{sous-groupes additifs stables par produit] et $1\in S_0$ [} sous-anneaux unitaires des $S_n$]
% archive sheet number: 9
\scanpage{11}
% 1–4
\uline{NB} \deleted{Le plus petit $A_*$} \inserted{a) b) c) impliquent} que $A_*$ est un sous-anneau cosimplicial de $S_*$. Le plus petit $A_*$ satisfaisant à a) b) c) est formé des

% 5–8
\[
\begin{aligned}
 A_n&=\mathrm{Im}\bigl(\mathbb Z[t_1,\ldots,t_n]\hookrightarrow k[[t_1,\ldots,t_n]]\bigr)\\
 &\simeq\Lambda[t_1,\ldots,t_n]
\end{aligned}
\]
où
\[
 \Lambda=\mathrm{Im}(\mathbb Z\unsure{-|\to} k)\quad
 \begin{cases}\simeq\mathbb Z&\text{ou}\\
 \simeq\mathbb Z/n\mathbb Z&\text{\deleted{pour} } n\in\mathbb N^*\ \text{\inserted{conven.}}\end{cases}
\]

% 9–13
Avec les conditions a) b) c) on construit la catégorie \deleted{$\unsure{\mathrm{Sfl}|\mathrm{Sf}}$} $\mathrm{Sfl}_{A_*}$ des \deleted{$A_*$}schémas \inserted{$A_*$-}formels (lisses). Dans cette catégorie, il y a des produits finis ($X_n\simeq(X_1)^n$), et des produits fibrés dans le cas ``transversal''.

% 13, left
\margin{Le foncteur ``oubli de $A_*$'' commute à ces lim $\leftarrow$}

% 14–18
\inserted{Soit $k_0=A_0\subset S_0=k$.} On obtient un foncteur $W_{A_*}$ de la cat. des $k_0$-modules libres de r.f. (et $=$ projectifs de type fini ?) vers les $A_*$-schémas formels, par
\[
 M_0\mapsto\mathrm{Spf}\,\widehat{\mathrm{Sym}}{}_k^*(M_{0k})
\]
(où $M_{0k}=M_0\otimes_{k_0}k$). C'est $=$ un foncteur à valeurs dans les objets groupes de $\mathrm{Sfl}_{A_*}$ i.e. vers les groupes $A_*$-formels.

% 14–18, left upper
% slanted
\margin{\deleted{\unsure{Notion} plus gén. $A_*$-schéma (fibré) sur $A_*\unsure{X|Y}$ et section sur $\unsure{A_X|A_Y}$, on a (suite \unsure{exacte|générale}) .\illegible{1}}}

% 14–18, left lower
% slanted
\margin{Prendre plus gén. des modules proj. de t. fini sur un $O_{A_*}(X)$ ..}

% 17–24, outer edge
% margin note written sideways
\margin{$O_{A_*}(X)\subset O(X)$. $\mathrm{Hom}_{A_*\text{ form}}(X,X_1)\subset\mathrm{Hom}_{k\text{-form}}(X,X_1)$ Modules \inserted{(lisses de t.f.)} $A_*$-formels sur les schémas $A_*$-formels .... \deleted{\unsure{Le|Le Sp}}\inserted{fibré} vectoriel $A_*$ formel associé $\overline V_{A_*}(M)$, ou $W_{A_*}(M)$}

% 19–25
Plus gén., si $\unsure{\mathrm{Lisc}|\mathrm{Linc}|\mathrm{Lin}}(k_0)$ est la catégorie des $k_0$-schémas augmentés \inserted{$k_0$-}isomorphes à \deleted{\unsure{Spec|Sp}} des \deleted{\unsure{E|k}} $\mathrm{Sp}\ k_0[\unsure{\mathrm{Sym}_k^*(M_{k_0})|\mathrm{Sym}^*(M_{k_0})|\widehat{\mathrm{Sym}}_k^*(M_0)}]$, \deleted{\illegible{1}}, il y a un foncteur $X\unsure{-|\to}\widehat X^{(A_*)}$ de $\mathrm{Lisc}(k_0)$ vers les $A_*$-schémas formels, commutant aux produits finis, et transformant objets groupes en objets groupes ....
% archive sheet number: 10
\scanpage{12}
% 1–2
Les conditions a) b) c) sont les \uline{conditions algébriques} élémentaires. On veut de plus

% 3, left
% slanted
\margin{après d) (condition différentielle)}

% 3–10
\deleted{\unsure{d}}\unsure{e)|d)} (Condition hensélienne) Si $u:X_n-X_n$ est un $A_*$-morphisme qui est un isomorphisme formel (i.e. si $f_1,\ldots,f_n\in A_n$ sont tels que $\varepsilon(f_1)\ldots\varepsilon(f_n)$ sont nilpotents dans $k$ et
\[
 \varepsilon\left(\det\left(\frac{\partial f_i}{\partial t_j}\right)\right)
 =\text{dét}\left(\varepsilon\left(\frac{\partial f_i}{\partial x_j}\right)\right)
 =\det\left(\frac{\partial f_i}{\partial\unsure{x_j|j|t_j}}(0)\right)
\]
est un élément inv. de $k$) alors $u^{-1}$ est aussi un $A_*$-morphisme ($\Leftrightarrow$ le foncteur des \deleted{$A_*$}schémas $A_*$ formels vers les schémas formels sur $k$ est conservatif).

% 11–16
\uline{Cor. 1}\qquad Si $f\in A_n$ et $f$ inversible dans $S_n$, i.e. $\varepsilon(f)=f(0)$ inversible dans $k$, alors $f^{-1}\in A_n$ ; \inserted{$\Longleftrightarrow$} \deleted{En particulier \inserted{\deleted{ou encore}}} si $f\in A_n$, $\varepsilon(f)=0$, alors $(1+f)^{-1}\in A_n$ [$\Longleftrightarrow$ si $f\in A_n$ et $\varepsilon(f)=1$, alors $f^{-1}\in A_n$ $\}\Longleftrightarrow A_n\cap\mathfrak m_n\subset\mathrm{rad}\,A_n$, où \deleted{$\mathfrak m_n=$} $\mathfrak m_n=(t_1,\ldots,t_n)S_n$ $\}$ \deleted{\unsure{au radical de $A_n\cap\mathfrak m_n$}}]

% 13, left
% slanted
\margin{moyennant condition d) (différentielle)}

% 17–20
En effet, considérons la \deleted{\unsure{multi}} l'endomorphisme
\[
 u:(x_1,\ldots,x_n,x_{n+1})\mapsto(x_1,\ldots,x_n,f(x_1,\ldots,x_n).x_{n+1})
\]
de $X_{n+1}$, il est un isom. formel d'inverse
\[
 (x_1,\ldots,x_n,x_{n+1})\mapsto(x_1,\ldots,x_n,f(x_1,\ldots,x_n)^{-1}x_{n+1}),
\]

% 21–24
donc comme \deleted{c'est}\inserted{$u$ est} un $A$-morphisme, par d) il en est de $=$ de $u^{-1}$, donc $g(t_1,\ldots,t_{n+1})=f(t_1,\ldots,t_n)^{\unsure{-1|-2}}t_{n+1}\in A_{n+1}$, \deleted{i.e. $X_{n+1}\to X_1$ \unsure{est un}} donc $g(t_1,\ldots,t_n,1)\in A_n$ i.e. $f^{-1}\in A_n$, cqfd.

% 22, left
% slanted
\margin{résulte de la condition différentielle}

% 24, left
% slanted
\margin{\unsure{incorrect|inconvénient}, $1$ n'est pas nilpotent ! Il faut \unsure{obliger|exiger} une propriété de divisibilité $t_{n+1}f\in A_{n+1}\Longrightarrow f\in A_{n+1}$}

% 25–27
\uline{Cor. 2}\qquad Soit $\inserted{P(Z)=}Z^N+f_1 Z^{N-1}+\ldots+f_N\in A_n[Z]$ polynôme unitaire, soit $z\in k_0$ tel que
\[
 P^\varepsilon(z)=z^N+\deleted{\unsure{f}}a_1 z^{N-1}\ldots+a_N=0,\quad
 \text{où }\deleted{\unsure{\xi}}\ a_i=\varepsilon(f_i)=f_i(0)\in k_0,
\]
% archive sheet number: 11
\scanpage{13}
% 1–7
et $(P^\varepsilon)'(z)=Nz^{N-1}+(N-1)a_1z^{N-2}\ldots+a_{N-1}$ soit $\in k^*$ (donc en fait $\in k_0^*$), de sorte qu'il existe un unique $F\in S_n$ telle que $P(F)=0$ et $\varepsilon(F)=z$. Sous ces conditions, on a $F\in A_n$. [En d'autres termes, le couple $(A_n,\mu_n=\mathfrak m_n\cap A_n)$ est (non seulement zariskien, selon le cor.~1, mais même) hensélien]

% 8–11
\deleted{En effet, \inserted{\unsure{quitte à remplacer|quitte à rempl.} $Z$ par $Z\unsure{+|-}z$, sps $z=0$, donc $\unsure{a_N|f_N}=\varepsilon(f_N)=0$} on regarde $X_{n+1}\to X_{n+1}$ donné par}
\[
\deleted{(x_1,\ldots,x_n,x_{n+1})\longmapsto (x_1,\ldots,x_n,x_{n+1}^{N+1}+f_1(x_1,\ldots,x_n)x_{n+1}^{N}\ldots)}
\]
\deleted{L'hyp. sur}

Plus généralement

% 12–16
\uline{Cor.~3} Soit $P\in A_{n+1}[t_1,\ldots,t_n,Z]$, $P=f_0+f_1Z+f_2Z^2+\ldots$ [les $f_i\in S_n(t_1,\ldots,t_n)$] supposons $\varepsilon(f_0)=0$ i.e. $\varepsilon(P)=0$, et $\varepsilon(f_1)\in k^*$ i.e. $\varepsilon(\frac{\partial P}{\partial Z})\in k^*$, de sorte qu'il existe une unique $F\in S_n$ telle que $P(t_1,\ldots,t_n,F(t_1,\ldots,t_n))=0$ et $\varepsilon(F)=0$. Alors $F\in A_n$.

% 17–24
Soit en effet $u:X_{n+1}\to X_{n+1}$ donné par $(x_1,\ldots,x_n,z)\longmapsto(x_1,\ldots,x_n,P(x_1,\ldots,x_n,z))$, alors l'hyp. dit que c'est compatible avec les ponctuations, et que c'est un isomorphisme formel. Comme l'hyp. $P\in A_{n+1}$ implique $u$ un $A_*$-morphisme, $u$ est un $A_*$ isomorphisme et par suite $u^{-1}$ est un $A_*$-morphisme. Donc
\[
\unsure{u^{-1*}|u^{-1}}(x_1,\ldots,x_n,z)=(x_1,\ldots,x_n,\deleted{\unsure{z}}\varphi(x_1,\ldots,x_n,z))
\]
avec $\varphi\in A_{n+1}(t_1,\ldots,t_n,t_{n+1})$, donc $\inserted{F=}\varphi(t_1,\ldots,t_n,0)\in A_n(t_1,\ldots,t_n)$

% 17–24, left
% margin note written sideways
\margin{a) b) c) d) impliquent-ils que $S_n$ soit fid. plat sur $A_n$ (si $k$ \deleted{\unsure{l'|le}}l'est \deleted{\unsure{plat|fini}} sur $k_0$) ?}

% 25–27
\deleted{d) (\unsure{Condition différentielle|Conditions différentielles}) Les $A_n$ sont \uline{stables} par les opérateurs différentiels standard}
\[
\deleted{D_{\unsure{p_i|p|p_j}}=D_{(p_1,\ldots,p_n)}\quad\left(=\frac{1}{p_1!\ldots p_n!}\frac{\partial^{|p|}}{\partial x_1^{p_1}\ldots\partial x_n^{p_n}}\right)}
\]
% archive sheet number: 12
\scanpage{14}
% 1–6
Soient $X$ un schéma $A_*$-formel et $Y$ un sous-schéma formel (lisse) de $X$. \deleted{On dit} On dit que $Y$ est un \deleted{\unsure{sous}} $A_*$-schéma $A_*$-formel s'il existe sur $Y$ une $A_*$ structure telle que l'inclusion can. soit un $A_*$-morphisme. Cette $A_*$-structure est alors unique.

% 6–14
En effet, il est immédiat qu'on peut (sans hyp. sur $Y$) trouver un \inserted{$A_*$} morphisme $X\to Y'$ ($Y'$ un schéma $A_*$-formel) tel que le morphisme induit $\deleted{X}Y\to Y'$ soit un isomorphisme formel. \deleted{\unsure{Considérons|Cons.}} Considérons l'isomorphisme inverse, et le composé
\diagram{\begin{tikzcd}[ampersand replacement=\&]
\& Y\arrow[dr] \& \\
Y'\arrow[ur,"\sim"]\arrow[rr] \& \& X
\end{tikzcd}}.
Utilisant (\unsure{e|c|d}), on voit que \deleted{$\unsure{A_*\text{ est}}$} $Y$ est un sous-schéma $A_*$-formel ssi $\deleted{\unsure{X\to}}Y'\to X$ est un $A_*$-morphisme ; et alors la $A_*$-structure de $Y$ est celle déduite.

% 15–17
\uline{Exemples} La diagonale de $X\times X$ est un sous $A_*$-schéma formel. Plus gén. celle de $X\times_Y X$ si $X\to Y$ est un $A_*$ morphisme submersif \deleted{(i.e. \unsure{formt. lisse|form. lisse|formel lisse})}

% 16–20, left
% slanted
\margin{Un morphisme formel $X\to Y$ entre \unsure{schs|objts} $A_*$-formels est un $A_*$-morphisme ssi son graphe dans $X\times Y$ est un sous-schéma $A_*$-formel}

% 18–27
Soit $\mathfrak X$ un schéma sur $\mathcal O_{A_*}(X)$ ($X$ un $A_*$-schéma formel) muni d'une section \inserted{$s_0$} sur $\mathcal O_{A_*}(X)$, supposons que $\mathfrak X$ soit lisse sur $\mathrm{Spec}\,\mathcal O_{A_*}(\deleted{X}X)$ aux points de cette section. Considérons le complété formel \inserted{$\widehat{\mathfrak X}\deleted{\unsure{^{s_0}}}$} de $\mathfrak X\otimes_{\mathcal O_{A_*}(X)}\mathcal O(X)$ le long de la section $s$ déduite de $s_0$. Alors une construction utilisant e) permet d'y mettre une \inserted{$A_*$-}structure, \inserted{telle que $X\to\widehat{\mathfrak X}\to X$ \unsure{soit|soient} \deleted{\unsure{de|des}} $A_*$-morphisme.} \deleted{\unsure{le|la}} $\mathfrak X$ \deleted{\unsure{est}} dépend $A_*$-fonctoriellement du triple $(X,\mathfrak X,s_0)$, sa formation est compatible avec les images inverses (par $X'\to X$\ldots)
% archive sheet number: 13
\scanpage{15}
% 1–4
d) (Condition différentielle) Tout $A_n$ est stable par les op. différentiels élémentaires $D_{\unsure{m|n|p}}$ [$\Longrightarrow$ Si $f(x_1,\ldots,x_n,y_1\ldots y_m)\in A_{n+m}$ est écrit sous la forme $\sum_{p\in\mathbb N^m}f_p(x)y^p$, alors les $f_p$ sont $\in A_n$]

% 1–5, left
% slanted
\margin{en car. 0 cela équivaut à la stabilité \inserted{des $A_n$} pour les $\frac{\partial}{\partial t_i}$.}

% 5–13
Ceci implique que les \unsure{algèbres|\deleted{algèbres}} $P^n_{X/k}$ sur $\mathcal O(X)$ ($X$ un \deleted{$A_*$} schéma $A_*$-formel) ont \deleted{une} $A_*$ structure canonique et fonctorielle, de $=$ $P^n_{X/Y}$ si $X\to Y$ est un $A_*$-morphisme submersif, de $=$ \deleted{\unsure{pour $\mathcal O$}} si $X\supset Y$ est un couple de schémas $A_*$-formels, $\mathcal O(Y)=\mathcal O(X)/J$, les invariants normaux $\mathcal O_X/J^n$ ont des $A_*$-structures canoniques, fonctorielles, compatibles aux images inverses transversales.

% 13–16
L'op. diff. universel d'ordre $n$\quad $\mathcal O(X)\longrightarrow P^n_{X/k}\deleted{\unsure{(X)|\text{etc.}}}$ \deleted{est un $A_*$ morphisme,} applique $\mathcal O_{A_*}(X)$ dans $(P^n_{X/k})_{A_*}\deleted{\unsure{^{X}|_{A_*}}}$ \ldots\ldots

% 16–21
\deleted{\unsure{D|Si}} les $\mathrm{Diff}_n(X/k)$ (\unsure{duales|duals|duaux} des $P^n_{X/k}$) ont \deleted{une} $A_*$-structure canonique, leur composition \deleted{est} laisse stables les parties $\mathrm{Diff}_n(X_{A_*}/k)$ des $A_*$-op. diff., et ces derniers stabilisent les \deleted{$\unsure{A(X)|A_*}$} $\mathcal O_{A_*}(X)$ \ldots\ Idem pour les op. diff. relatifs, pour $X\to Y$ $A_*$-morphisme \deleted{\unsure{rel. lisse|lisse}}\inserted{submersif}.

% 22–25
\uline{Cor} $A_n\subset k_0[[t_1,\ldots,t_n]]$. \deleted{\unsure{Alors}} D'où un foncteur $\unsure{M\text{-ent}|M_*|\text{allant}}$ \deleted{$\int$ \unsure{(système, catégorie)|catégorie}} des schémas $A_*$-formels vers les schémas \deleted{Soit $X\to Y$ un morphisme submersif de schémas formels, sur $X$} formels sur $k_0$ \ldots
% archive sheet number: 14
\scanpage{16}
% 1–7
Soit $p:X\to Y$ un morphisme submersif de schémas formels (lisses), \unsure{supposons} $X$ muni d'une $A_*$-structure. On dit que $Y$ admet une $A_*$-structure quotient de celle de $X$ (moyennant $p$) s'il admet une \inserted{$A_*$} structure telle que $p$ devienne un $A_*$-morphisme. Celle-ci est encore unique : en effet, il existe (sans condition sur $Y$)

% 8–14
un \deleted{\unsure{sous $A$|sous-$A_*$}} sous-schéma $A_*$-formel $Y'$ de $X$, tel que le morphisme formel induit $Y'\to Y$ soit un isomorphisme. \inserted{Considérons} le composé
\diagram{\begin{tikzcd}[ampersand replacement=\&]
\& Y\arrow[dr,"\sim"] \& \\
X\arrow[ur,"p"]\arrow[rr] \& \& Y'
\end{tikzcd}}
\deleted{\unsure{alors}}, on voit alors grâce à c) que la $A_*$-structure quotient existe ssi $X'\to Y'$ est un $A_*$-morphisme, et la $A_*$-structure sur $Y$ est celle déduite de celle de $Y'$.

% 15–19
\deleted{\unsure{Une co|Une ca}} Soit $X\xrightarrow{p}Y$ un morphisme formel submersif. Alors $R=X\mathbin{\widehat{\times}}_Y X\subset X\mathbin{\widehat{\times}}X$ est un \inserted{sous-schéma formel lisse, et c'est} relation d'équivalence [stable par symétrie, contient diagonale, transitif\ldots]. \deleted{\unsure{Inversement}}\inserted{Plus précisément},

% 20–25
$X$ étant donné, il y a corr. biunivoque entre ces quotients formels $Y$, et les $R$ rel. d'équiv. formelles. Si $\deleted{A}X$ est muni d'une $A_*$-structure, alors si $Y$ admet une $A_*$-structure quotient, on voit que $R$ est un sous-schéma $A_*$-formel de $X\mathbin{\widehat{\times}}X$ i.e. il a une $A_*$-structure induite.

% 26–28
f) (Condition intégrale) La réciproque est vraie, i.e. si $R$ est une relation d'équival. formelle dans $X$, qui soit une $A_*$-sous-structure de $X\mathbin{\widehat{\times}}X$, alors
% archive sheet number: 15
\scanpage{17}
% 1–2
le quotient \inserted{formel} $Y=X/R$ admet une $A_*$-structure quotient.

% 3–11
[Heuristiquement, \unsure{étant|ét.} donné un $A_*$-feuilletage de $X$, la variété formelle des feuilles admet une $A_*$-structure quotient. En car. 0, il est ``bien connu'' que les $A_*$-feuilletages correspondent aux ``systèmes de Pfaff complètement ``intégrables'' \unsure{sur|\deleted{sur}} $X$, et la ``condition intégrale'' signifie que l'intégration de ce système (\unsure{effectué|effectuée}, disons, par le morphisme formel plus haut $X\to Y'$) se fait par un $A_*$-morphisme\ldots.
% archive sheet number: 16
\scanpage{19}
% 1–5
\uline{Th. de préparation ``formel''} Soit $A$ un anneau linéairement topologisé, séparé et complet, $Z$ une indéterminée,
\[
F\deleted{(\unsure{Z})}=\sum_{i\geq 0}a_i Z^i\in A[[Z]],\quad N\in\mathbb N\ \text{tels que}
\]
(i)\qquad $a_i$ est top.$^{\mathrm t}$ nilpotent pour $0\leq i\leq N-1$

(ii)\qquad $a_N$ est inversible

% 6–10
Alors $F$ est un élément régulier de $A[[Z]]$ et $A[[Z]]/F$ a comme base $(Z^i)_{0\leq i\leq N-1}$, en d'autres termes tout élément $G$ de $A[[Z]]$ s'écrit de façon unique sous la forme
\[
\begin{aligned}
G&=QF+\sum_{0\leq i\leq N-1}b_iZ^i & Q&\in A[[Z]]\\
&&b_i&\in A
\end{aligned}
\]

% 11–15
\uline{Démonstration} Comme $A$ est séparé et complet pour sa top. linéaire, on est ramené aussitôt au cas où $A$ est discret, alors (i) signifie que les $a_i$ $(0\leq i\leq N-1)$ sont top. nilpotents. Soit $\mathfrak m$ l'idéal \inserted{(nilpotent)} qu'ils engendrent, $A_0=A/\mathfrak m$, et désignons par $\dot a_i,\dot F$ etc les réductions mod $\mathfrak m$.

% 16–22
\deleted{\unsure{Alors}} Donc on a \deleted{$\unsure{\dot a}$}
\[
\dot a_i=0\quad\text{pour}\quad 0\leq i\leq N-1
\]
i.e
\[
\begin{aligned}
\dot F&=\dot a_{N\deleted{\unsure{+1}}}Z^{N\deleted{\unsure{+1}}}+\ldots\\
&=\dot a_{N\deleted{+1}}Z^{N\deleted{+1}}(1+\ldots)
\end{aligned}
\]
donc, comme $\dot a_{N\deleted{\unsure{+1}}}$ est inversible par (ii), on voit que $\dot F$ est de la forme
\[
\dot F=\dot U Z^N,
\]
\deleted{\unsure{donc les|de sorte que les}} donc
\[
\unsure{\dot A_0|A_0}[[Z]]/\dot F\simeq A_0[[Z]]/Z^N\deleted{{}=A_0[Z]/Z^{\unsure{n|N}}}
\]
a comme base les $Z^i$ $(0\leq i\leq N-1)$

% 23–31
\deleted{i.e. (i) les $Z^i$ $(0\leq i\leq N-1)$ engendrent $A[[Z]]/FA[[Z]]$}

\deleted{En effet, ils engendrent mod $\mathfrak m$, on gagne donc par Nakayama.}

\deleted{\deleted{\unsure{Pour ce qui précède|Pour ce point il suffit de}}}

\deleted{(ii) $F$ est régulier. En effet, \deleted{\unsure{par récurrence sur l'ordre}} \inserted{le $\unsure{A\text{-module}|A_0\text{-module}}$ $A[[Z]]$ filtré $\mathfrak m$-adiquement} a un gr. associé isomorphe à \deleted{$\mathfrak m^n/\mathfrak m$} $\mathrm{Gr}^{*}(A)^{\mathbb N}$. \deleted{de nilpotence de $\mathfrak m$, SPS $\mathfrak m^2=0$. Soit $Q\in A[[Z]]/$ t.q. que $QF=0$, donc $Q\cdot$} et l'opération de multiplication \deleted{de $F$} par $F$ dessus, sur le gradué, correspond à : l'op. de multiplication par $\dot F$, qui est injective, OK.}
% archive sheet number: 18
\scanpage{20}
% 1–4
\deleted{(iii) Etude du cas où $\deleted{\unsure{F=|F\text{ est}}}\left\{\begin{aligned}a_i&=0\quad\text{pour }i\geq N+1\\ a_N&=1\end{aligned}\right.$}

\deleted{i.e. du cas où}
\[
\deleted{F=Z^N+a_{N-1}Z^{N\unsure{-|+}1}\ldots+a_0}
\]
\deleted{les $a_i$ $(0\leq i\leq N-1)$ nilpotents.}

% 5–12
\[
\deleted{\underset{\substack{\text{base }(Z^i)_{0\leq i\leq N-1}\\\text{sur }A}}{A[Z]/FA[Z]}\longrightarrow A[[Z]]/FA[[Z]]}
\]
\deleted{est alors surjectif (car les $Z^i$ $(0\leq i\leq N-1)$ engendrent $A[[Z]]/FA[[Z]]$ comme $A$-module) pr dire que c'est bijectif i.e. que $(FA[[Z]])\cap A[Z]=FA[Z]$ (il revient au $=$ de dire que les \deleted{\unsure{deux classes|classes de}} $Z^i$ $(0\leq i\leq N-1)$ forment une base de $A[[Z]]/FA[[Z]]$).}

% 13–14
\deleted{Soit donc $a_0+a_1Z\ldots+a_{N-1}Z^{N-1}$ qui soit multiple de $F$ dans $A[[Z]]$, \unsure{soi.|soit|i.e.}}

% 15–17
Considérons l'application
\[
\begin{aligned}
\varphi:\quad A^N+A[[Z]]&\longrightarrow A[[Z]]\\
\big((b_i)_{0\leq i\leq N-1},Q\big)&\longrightarrow\sum_{0\leq i\leq N-1}\deleted{\unsure{Q_i|Q}}b_iZ^i+QF
\end{aligned}
\]

% 18–23
l'application déduite par passage aux gradués associés pour la filtration $\mathfrak m$-adique (finie) est
\[
(\mathrm{Gr}\,A)^N\times(\mathrm{Gr}\,A)[[Z]]\longrightarrow(\underline{\unsure{\mathrm{Gr}'|\mathrm{Gr}}\,A})[[Z]]
\]
\unsure{donné|donnée} \deleted{par} en degré $n$ par
\[
((b_i),Q)\longrightarrow\sum_{0\leq i\leq N-1}b_iZ^i+Q\dot F
\]

% 24–26
et comme $\dot F=\deleted{\dot a_{N+1}}\alpha Z^N$ $(\alpha\unsure{=|-}a_{N+1}\in A_0^*)$ ceci est manifestement bijectif, donc $\varphi$ aussi, cqfd.
% archive sheet number: 19
